\begin{lemma}[Lemma 2.3]
Let $n\ge2$, $\gamma>0$ and $\lambda_0,\lambda_1\in(0,1)$, and suppose the Sobolev--Poincar\'e inequality on balls holds in dimension $n$
(\noderef{CitedSobolevPoincareBall}).  Then there exists
$\theta_0=\theta_0(n,\gamma,\lambda_0,\lambda_1)>0$ with the following property.
Let $x_0\in\R^n$, let $R>0$, let $f\in W^{1,2}(B_R(x_0))$ with weak gradient $g$
(\noderef{MemW1pOn}), suppose $f$ is essentially bounded above on $B_R(x_0)$ and
put $M=\esssup_{B_R(x_0)}f$.  Suppose $k<M$ and that the De Giorgi level bound
(2.1) of the thesis with constant $\gamma$ (\noderef{DeGiorgiLevelBound}) holds for
every level $\ell\in[k,\lambda_0k+(1-\lambda_0)M]$ and all radii
$\rho_0,\rho_1\in[\lambda_1R,R]$ with $\rho_0<\rho_1$.  If
\[
\bigl|\{x\in B_R(x_0):f(x)>k\}\bigr|\;\le\;\theta_0\,|B_R| ,
\]
then
\[
f\le\lambda_0k+(1-\lambda_0)M\qquad\text{a.e.\ in }B_{\lambda_1R}(x_0).
\]

The hypothesis \noderef{CitedSobolevPoincareBall} is the cited result (C6) of the
paper's Note for the formalization, carried here as an explicit argument because
the proof uses it --- it enters through the constant $\beta(n)$ of \noderef{DeGiorgiLevelSetMeasure} --- which is what the
Note permits for the results the thesis quotes without proof.  Although it is
bound before the existential for $\theta_0$, it introduces no dependence of
$\theta_0$ on any quantity outside the list the paper names: (C6) is a statement
about $n$ alone, and $n$ is already bound first, so the hypothesis carries no
numerical content of its own.
\end{lemma}
\begin{proof}
Write $\omega_n=|B_1|$, so $|B_\sigma|=\omega_n\sigma^n$, let $\beta(n)$ be the
constant of \noderef{DeGiorgiLevelSetMeasure} --- available because the
hypothesis \noderef{CitedSobolevPoincareBall} of the present lemma is exactly the
hypothesis that node's statement requires --- and set
\[
q=\tfrac n{n-1},\quad \varepsilon=\tfrac1{n-1},\quad
b=\bigl[(1-\lambda_0)(1-\lambda_1)\bigr]^{\frac n{1-n}},\quad
C_1=\frac{2\beta(n)\sqrt\gamma}{\omega_n\lambda_1^{\,n+1}},\quad
C_2=\frac{C_1}{\lambda_0(1-\lambda_0)(1-\lambda_1)},\quad
c=C_2^{\,q} .
\]
All of these depend only on $n,\gamma,\lambda_0,\lambda_1$.  Note $\varepsilon>0$
and $q=1+\varepsilon$ because $n\ge2$, and $b>1$ because
$(1-\lambda_0)(1-\lambda_1)\in(0,1)$ while $\frac n{1-n}<0$.  Define
\[
\theta_1=\frac{\lambda_1^{\,n}}2,\qquad
\theta_2=\frac{c^{-1/\varepsilon}\,b^{-1/\varepsilon^2}}{\omega_n},\qquad
\theta_0=\min\{\theta_1,\theta_2\} .
\]
Then $\theta_0>0$ and $\theta_0$ depends only on $n,\gamma,\lambda_0,\lambda_1$;
note $\theta_1=|B_{\lambda_1R}|/(2|B_R|)$ for every $R>0$, which is why it carries
no $R$.  Fix $x_0,R,k,M,f,g$ as in the statement and assume
$|A^+_{k,R}|\le\theta_0|B_R|$, where throughout
$A^+_{\sigma,\rho}=\{x\in B_\rho(x_0):f(x)>\sigma\}$.

\emph{Step 1: the two sequences.}  For $h=0,1,2,\dots$ put
\[
r_h=\lambda_1R+(1-\lambda_1)^{h+1}R,\qquad
k_h=k+(1-\lambda_0)(M-k)-(1-\lambda_0)^{h+1}(M-k).
\]
Then $r_0=\lambda_1R+(1-\lambda_1)R=R$; $r_h$ is strictly decreasing with
$r_h\downarrow\lambda_1R$, so $\lambda_1R<r_h\le R$ for every $h$; and
\[
r_h-r_{h+1}=(1-\lambda_1)^{h+1}R-(1-\lambda_1)^{h+2}R=\lambda_1(1-\lambda_1)^{h+1}R .
\]
Also $k_0=k$; $k_h$ is strictly increasing with
$k_h\uparrow k+(1-\lambda_0)(M-k)=\lambda_0k+(1-\lambda_0)M$, so
$k\le k_h<\lambda_0k+(1-\lambda_0)M$ for every $h$; and
\[
k_{h+1}-k_h=(1-\lambda_0)^{h+1}(M-k)-(1-\lambda_0)^{h+2}(M-k)
=\lambda_0(1-\lambda_0)^{h+1}(M-k) ,
\]
which is positive since $M>k$.  Finally put
\[
\nu_h=\frac{|A^+_{k_h,r_h}|}{R^n}\ \ge 0 .
\]

\emph{Step 2: the level-set estimate at stage $h$.}  Fix $h$.  We apply
\noderef{DeGiorgiLevelSetMeasure} on the ball $B_{r_{h+1}}(x_0)$ at the levels
$k_h\le k_{h+1}$.  Its hypotheses hold: $r_{h+1}>0$; $f\in W^{1,1}(B_{r_{h+1}})$
with weak gradient $g$, because $f\in W^{1,2}(B_R)$ with weak gradient $g$
(\noderef{MemW1pOn}), $B_{r_{h+1}}\subseteq B_R$ has finite measure, so
H\"older's inequality turns the $L^2$ bounds on $f$ and $|g|$ into $L^1$ bounds,
while the weak-gradient identity restricts to test functions supported in
$B_{r_{h+1}}$ because such a test function vanishes off $B_{r_{h+1}}$; and the
nondegeneracy hypothesis holds by the following computation, which is where
$\theta_0\le\theta_1$ is consumed.  Since $k_h\ge k$ and $r_{h+1}\le R$ we have
$A^+_{k_h,r_{h+1}}\subseteq A^+_{k,R}$, so
\[
|A^+_{k_h,r_{h+1}}|\le|A^+_{k,R}|\le\theta_0|B_R|\le\theta_1|B_R|
=\tfrac12|B_{\lambda_1R}| ,
\]
and since $B_{\lambda_1R}\subseteq B_{r_{h+1}}$,
\[
\bigl|B_{r_{h+1}}\setminus A^+_{k_h,r_{h+1}}\bigr|
\ge\bigl|B_{\lambda_1R}\bigr|-\bigl|A^+_{k_h,r_{h+1}}\bigr|
\ge\tfrac12\bigl|B_{\lambda_1R}\bigr|>0 .
\tag{2.2a}
\]
\noderef{DeGiorgiLevelSetMeasure} therefore gives
\[
(k_{h+1}-k_h)\,\bigl|A^+_{k_{h+1},r_{h+1}}\bigr|^{1-\frac1n}
\le\beta(n)\,\frac{r_{h+1}^{\,n}}{\bigl|B_{r_{h+1}}\setminus A^+_{k_h,r_{h+1}}\bigr|}
\int_{A^+_{k_h,r_{h+1}}\setminus A^+_{k_{h+1},r_{h+1}}}|g|\,dx .
\tag{2.2}
\]
Using (2.2a), $r_{h+1}\le R$, $|B_{\lambda_1R}|=\omega_n\lambda_1^nR^n$ and
enlarging the domain of integration to $A^+_{k_h,r_{h+1}}$,
\[
\lambda_0(1-\lambda_0)^{h+1}(M-k)\,\bigl|A^+_{k_{h+1},r_{h+1}}\bigr|^{1-\frac1n}
\le\frac{2\beta(n)}{\omega_n\lambda_1^{\,n}}\int_{A^+_{k_h,r_{h+1}}}|g|\,dx .
\tag{2.3}
\]

\emph{Step 3: estimating the gradient integral by (2.1) of the thesis.}  Apply
\noderef{DeGiorgiLevelBound} with constant $\gamma$ at the level $\ell=k_h$ and
the radii $\rho_0=r_{h+1}$, $\rho_1=r_h$.  This is permitted by the hypothesis of
the lemma: $k_h\in[k,\lambda_0k+(1-\lambda_0)M]$ by Step 1, and
$r_{h+1},r_h\in[\lambda_1R,R]$ with $r_{h+1}<r_h$, also by Step 1.  Its second
clause reads
\[
\int_{B_{r_{h+1}}\cap\{f>k_h\}}|g|^2dx
\le\frac{\gamma}{(r_h-r_{h+1})^2}
\Bigl(\esssup_{B_{r_h}}(f-k_h)^+\Bigr)^{2}\bigl|A^+_{k_h,r_h}\bigr| .
\]
Now $B_{r_h}\subseteq B_R$ and $M=\esssup_{B_R}f$, so $f\le M$ a.e.\ on
$B_{r_h}$; hence $(f-k_h)^+\le(M-k_h)^+=M-k_h\le M-k$ a.e.\ on $B_{r_h}$ (using
$k\le k_h<M$), and therefore
$\esssup_{B_{r_h}}(f-k_h)^+\le M-k$.  Note this uses the essential-boundedness
hypothesis on $f$ only through the inequality $f\le M$ a.e.; no supremum is
assumed attained.  With $r_h-r_{h+1}=\lambda_1(1-\lambda_1)^{h+1}R$ we get
\[
\int_{A^+_{k_h,r_{h+1}}}|g|^2dx
\le\frac{\gamma\,(M-k)^2}{\lambda_1^2(1-\lambda_1)^{2h+2}R^2}\bigl|A^+_{k_h,r_h}\bigr| .
\]
By the Cauchy--Schwarz inequality on $A^+_{k_h,r_{h+1}}$, and then
$|A^+_{k_h,r_{h+1}}|\le|A^+_{k_h,r_h}|$ (valid since $r_{h+1}<r_h$),
\[
\int_{A^+_{k_h,r_{h+1}}}|g|\,dx
\le\bigl|A^+_{k_h,r_{h+1}}\bigr|^{1/2}
\Bigl(\int_{A^+_{k_h,r_{h+1}}}|g|^2dx\Bigr)^{1/2}
\le\frac{\sqrt\gamma\,(M-k)}{\lambda_1(1-\lambda_1)^{h+1}R}\bigl|A^+_{k_h,r_h}\bigr| .
\]
Substituting into (2.3),
\[
\lambda_0(1-\lambda_0)^{h+1}(M-k)\,\bigl|A^+_{k_{h+1},r_{h+1}}\bigr|^{1-\frac1n}
\le\frac{2\beta(n)\sqrt\gamma}{\omega_n\lambda_1^{\,n+1}}\,
\frac{M-k}{(1-\lambda_1)^{h+1}R}\bigl|A^+_{k_h,r_h}\bigr| ,
\tag{2.4}
\]
which is (2.4) of the thesis of the paper in the form we need.

\emph{Step 4: the recursion.}  Divide (2.4) by
$\lambda_0(1-\lambda_0)^{h+1}(M-k)>0$:
\[
\bigl|A^+_{k_{h+1},r_{h+1}}\bigr|^{1-\frac1n}
\le\frac{C_1}{\lambda_0(1-\lambda_0)^{h+1}(1-\lambda_1)^{h+1}}
\cdot\frac{\bigl|A^+_{k_h,r_h}\bigr|}{R}
\le\frac{C_2}{\bigl[(1-\lambda_0)(1-\lambda_1)\bigr]^{h}}
\cdot\frac{\bigl|A^+_{k_h,r_h}\bigr|}{R} ,
\tag{2.5}
\]
by the definitions of $C_1$ and $C_2$.  Raise (2.5) to the power
$q=\frac n{n-1}>0$ and divide by $R^n$.  Writing
$|A^+_{k_h,r_h}|=R^n\nu_h$ and $|A^+_{k_{h+1},r_{h+1}}|=R^n\nu_{h+1}$, the
left-hand side becomes
$R^{n}\nu_{h+1}\big/R^{n}=\nu_{h+1}$ after division, because
$\bigl(|A^+_{k_{h+1},r_{h+1}}|^{1-\frac1n}\bigr)^{q}=|A^+_{k_{h+1},r_{h+1}}|$,
and the right-hand side becomes
\[
\frac{C_2^{\,q}}{\bigl[(1-\lambda_0)(1-\lambda_1)\bigr]^{hq}}
\cdot\frac{\bigl(R^n\nu_h\bigr)^{q}}{R^{q}\,R^n}
=c\,b^{\,h}\,\nu_h^{\,q},
\]
since the power of $R$ is $nq-q-n=\frac{n^2}{n-1}-\frac n{n-1}-n
=\frac{n(n-1)}{n-1}-n=0$, and
$\bigl[(1-\lambda_0)(1-\lambda_1)\bigr]^{-hq}=b^{\,h}$ by the definition of $b$.
As $q=1+\varepsilon$ we have obtained
\[
\nu_{h+1}\le c\,b^{\,h}\,\nu_h^{\,1+\varepsilon}\qquad(h=0,1,2,\dots),
\tag{2.6}
\]
the recursion in the $b^{\,h}$ form.

\emph{Step 5: the De Giorgi iteration.}  We check the hypotheses of
\noderef{DeGiorgiSequenceLemma} for the sequence $(\nu_h)$ with the constants
$c,\varepsilon,b$ above: $c>0$, $\varepsilon>0$, $b>1$, $\nu_h\ge0$ for all $h$,
and (2.6) is the required recursion.  It remains to verify the threshold
condition $\nu_0\le c^{-1/\varepsilon}b^{-1/\varepsilon^2}$, and this is where
$\theta_0\le\theta_2$ is consumed:
\[
\nu_0=\frac{\bigl|A^+_{k_0,r_0}\bigr|}{R^n}=\frac{\bigl|A^+_{k,R}\bigr|}{R^n}
\le\frac{\theta_0|B_R|}{R^n}=\theta_0\,\omega_n\le\theta_2\,\omega_n
=c^{-1/\varepsilon}b^{-1/\varepsilon^2} ,
\]
using $k_0=k$, $r_0=R$ from Step 1 and $|B_R|=\omega_nR^n$.  Hence
\noderef{DeGiorgiSequenceLemma} applies and gives
$\lim_{h\to\infty}\nu_h=0$.

\emph{Step 6: conclusion.}  Put $k_\infty=\lambda_0k+(1-\lambda_0)M$.  For every
$h$ we have $k_h\le k_\infty$ and $\lambda_1R\le r_h$, hence
\[
\bigl\{x\in B_{\lambda_1R}:f(x)>k_\infty\bigr\}
\subseteq\bigl\{x\in B_{r_h}:f(x)>k_h\bigr\}=A^+_{k_h,r_h} ,
\]
so $\bigl|\{x\in B_{\lambda_1R}:f>k_\infty\}\bigr|\le|A^+_{k_h,r_h}|=R^n\nu_h$.
Letting $h\to\infty$ and using Step 5, the left-hand side is $\le0$, hence
\[
\bigl|\{x\in B_{\lambda_1R}(x_0):f(x)>\lambda_0k+(1-\lambda_0)M\}\bigr|=0 ,
\]
that is, $f\le\lambda_0k+(1-\lambda_0)M$ almost everywhere in
$B_{\lambda_1R}(x_0)$, which is the assertion.
\end{proof}
